\chapter{96}

\section{Bettmeralp/VS (CH), Dienstag, 11.
September}



Ava stellte gerade eine Kaffeetasse unter den Ausguss der Kaffeemaschine, als
sie Louis vom Stall über den Flur in die Küche kommen hörte.

Erneut hatte es geklappt, dass sie sich weder beim Einschlafen noch beim
Aufwachen direkt begegnet waren. Louis hatte zwar kaum geschlafen. Doch
intensiv mit sich selbst beschäftigt, hatte er Avas Anwesenheit
ausgeblendet.

Ava zuckte zusammen und sah ihn auf sie zukommen. Auch für Louis war die
Begegnung beklemmend. Er hatte versucht, sich darauf vorzubereiten. Nach
diesem verheissungsvollen Tag gestern. Es war einer wahren Entblössung
gleich gekommen. Und an Schmach nicht zu übertreffen gewesen. Und doch
war dieses nicht mehr auszuhaltende Geheimnis endlich, endlich gelüftet.
Joh hatte die andern, wie vereinbart, informiert. Ciro war weg.
Wenigstens hier oben. Wie es mit Louis weitergehen würde, stand in den
Sternen. Von nun an gab es ihn immerhin nicht mehr doppelt.

Ava gab dem Drang nach, ihn zu erlösen.

«Louis. Wusste ich’s doch.»

Sie versuchte, verständnisvoll zu klingen.

Louis fand keine Worte. Er nahm sich eine Tasse aus dem Regal und
stellte sich neben sie. 
Ava blieb, die Tasse mit beiden Händen umklammert, neben der Maschine
stehen. Wieder dieser betörende Duft. Wenigstens, dachte sie, 
ist etwas Echtes von ihm übrig geblieben.
Sie schaute zu, wie sich Louis’ Tasse füllte, als sähe sie zum
ersten Mal Kaffee fliessen.

Louis stellte die Maschine aus. Ava wartete. Ein Knistern lag in der
Luft. Louis kämpfte. Es gab kein Weglaufen mehr. Auch nicht vor ihr.

Er gab sich einen Ruck und schaute sie direkt an.

«Es ist schwer für mich, Ava. Wenn, wenn du willst, erklär ich dir, wer
ich wirklich bin.»

Sie nickte bereits bei Louis’ ersten Worten. Ihn zu erlösen, kam in
dieser beklemmenden Situation ihrer eigenen Erlösung gleich.

«Ja,» sie suchte nach seinem Namen, «ja, Louis, das ist gut.»

«Setzen wir uns draussen auf die Bank?»

Mit einer Handbewegung liess Louis Ava vorangehen.

\sectionbreak

Es war ganz selbstverständlich gekommen. Ava hatte sich nicht mehr
zurückhalten können. Louis hatte zu traurig ausgesehen, weinte
plötzlich, leise, und schliesslich herzzerreissend. Es hatte mit einer
kleinen Berührung seiner Hand begonnen. Die Stunden gehörten einzig
Louis’ Kummer und Avas Anteilnahme. Und während sie nach und nach näher
an ihn heranrückte, zogen sich Zeit und Distanz diskret zurück und sie
hielt ihn im Arm. Und dass in diesem Moment Manuela von ihrem Einkauf
zurückkam, sie beide in dieser Situation sah, ein Ohlala-Gesicht machte
und taktvoll verschwand, war Ava wenig peinlich.

Louis löste sich aus der Umarmung. Er sah anders aus als sie ihn bisher
wahrgenommen hatte. Weicher vielleicht, ehrlicher, dachte sie. Er war
nun wohl sich selbst, Louis, und nicht dieser erfundene Ciro, dem sie
mit der Zeit misstraut hatte. Und wenn sie ganz ehrlich war, hatte er
sie auch immer wieder genervt.

«Phua, Ci…eh Louis.»

Ava atmete durch. Die Kapellenuhr oben schlug zwölf. Sie hatten die Zeit
vergessen. Louis hatte ewig erzählt. Sie hatte gebannt zugehört,
manchmal Zwischenfragen gestellt. Vor allem aber hatte sie gestaunt.
Sein Gesichtsausdruck hatte sich von Minute zu Minute verwandelt. Wie er
sie anschaute und den Blick hielt. Wie fokussiert er war.

Louis setzte sich gerade auf. Er wischte sich über die Stirn, griff nach
seiner Tasse und leerte den längst kalt gewordenen Kaffee in sich
hinein.

«Ich, ich fass es nicht, Ava.»

Ava stutzte.

«Was ich dir alles sage, mein ich. Ich, ehm, so viel habe ich bis heute
kaum jemandem von mir erzählt.»

Louis hob kurz die Arme und liess sie schliesslich neben sich hängen. Er
sah entrückt aus.

Sie sassen stumm nebeneinander.

Über ihnen, verteilt über das weite Hausdach, begann ein zaghaftes
Klopfen. Sie schauten gleichzeitig hoch. Louis schmunzelte. Erste
Regentropfen. Endlich, nach all den dürren Tagen.

Ava atmete hörbar ein.

«Kennst du Ultimo?»

Louis drehte sich ganz zu ihr um.

«Jaa,» er zog das Wort in die Breite und kniff die Augen zusammen,
«schon gehört.»

«Seine Songs gefallen mir. Daniel und ich, wir, wir haben uns seine
Lieder oft angehört.»

Ava machte eine Pause. Mit einem Mal fühlte sie sich befangen. Louis
legte den Kopf schräg und suchte ihren Blick.

«Also, weisst du, eh, Louis, als du erzähltest, da, da fiel mir immer
wieder ein Song von ihm ein.»

Louis’ Gesichtsausdruck verriet Interesse. Er streckte den Kopf vor und
nickte erwartungsvoll.

«Ich weiss nicht, ob es passt,» sie klang vorsichtig, «ich denk, der
Text von \emph{«Sogni Appesi»} könnte dir gefallen. Der Song handelt
von ungelebten Träumen, von schweren Zeiten und…»

Sie stockte. Louis legte seine Hand auf ihren Arm. Die Berührung fühlte
sich entlastend an. Jetzt lachte sie. Verrückt, wie offen sie sich
gleich wieder fühlte.

«Wollen wir ihn uns anhören?»

«Gerne, Ava, machen wir,» Louis zog sein Handy aus der Hosentasche,
schaute es für einen Moment gebannt an, dachte, dass das Gerät seit
gestern seine Gefährlichkeit hier oben verloren hatte, und fuhr fort,
«ich hol den Boomer raus.»

Beim Hochgehen und auf oberster Stufe, blieb Louis stehen. Er drehte
sich um.

«Danke.»

Avas Augen schlossen sich. Dann nickte sie. Sie fühlte sich verbunden.
Mit sich, mit Louis, mit der Gutmütigkeit dieses Ortes, mit dem
regelmässigen Klackern der Regentropfen, vielmehr mit jedem Regentropfen
einzeln, und gerade stimmte alles.

Über das offene Küchenfenster erreichte sie der Duft gebratener
Kartoffeln.

\qrblock{Ultimo \emph{«Sogni Appesi»}}{media/QR-Codes/050_Sogni_Appesi_Ultimo_Spotify.png}
